\section{Introduction}

\subsection{Task Description}

Multimodal Large Language Models (MLLMs) have demonstrated powerful capabilities
in visual question answering and image-text reasoning,
yet they still suffer from hallucinations,
generating outputs that are inconsistent with visual inputs,
contradict world knowledge,
or contain flawed reasoning.
This experiment focuses on hallucination detection and evaluation in MLLMs.
The assessment requirements are as follows:

\begin{enumerate}
  \item Hallucination definition and classification:
    define at least two types of hallucinations,
    with no fewer than two examples per type and clear explanations.
  \item Automated detection methods and evaluation metrics:
    construct at least two detection methods (e.g., rule-based, MLLM Judge)
    and design at least two evaluation metrics.
  \item Experimental evaluation (main task):
    evaluate at least two models
    (including one closed-source and one open-source)
    on general visual question answering
    and mathematical image-text reasoning datasets.
  \item Human alignment verification:
    construct 20--50 human-annotated samples,
    compare automated detection results with human judgments,
    and analyze inconsistent cases.
  \item Analysis requirements:
    include analysis of CoT's impact on hallucinations,
    cross-model comparison,
    and in-depth analysis of at least two typical failure cases.
  \item Bonus:
    extend to other image-text tasks, specific domains,
    or conduct cross-task generalization analysis.
\end{enumerate}

\subsection{Evaluated Models}

This experiment selects four representative MLLMs,
covering both closed-source and open-source models
with diverse parameter scales and architectural designs,
to provide a multi-dimensional comparative perspective.
\Cref{tab:model_config} summarizes the key configurations.

\begin{table}[htbp]
  \centering
  \caption{Model configurations.}
  \label{tab:model_config}
  \begin{tabular}{lllll}
    \toprule
    Model & Type & Parameters & API Protocol & Platform \\
    \midrule
    GPT-5.4-mini       & Closed-source   & Undisclosed          & OpenAI Responses & right.codes \\
    Gemini 2.5 Flash   & Closed-source   & Undisclosed          & OpenAI Chat      & right.codes \\
    Qwen3.5-35B-A3B    & Open-source MoE & 35B (3B activated)   & OpenAI Chat      & SiliconFlow \\
    Qwen3-VL-235B-A22B & Open-source MoE & 235B (22B activated) & OpenAI Chat      & SiliconFlow \\
    \bottomrule
  \end{tabular}
\end{table}

\subsection{Evaluation Datasets}

Three standard benchmarks are selected,
covering object perception, mathematical reasoning,
and medical VQA tasks respectively.
All datasets are sourced from Hugging
Face.\footnote{POPE: \url{https://huggingface.co/datasets/lmms-lab/POPE};
MathVista: \url{https://huggingface.co/datasets/AI4Math/MathVista};
VQA-RAD: \url{https://huggingface.co/datasets/flaviagiammarino/vqa-rad}.}
\Cref{tab:dataset_overview} provides a summary;
detailed dataset descriptions are given in \Cref{sec:appendix}.

\begin{table}[htbp]
  \centering
  \caption{Dataset overview.}
  \label{tab:dataset_overview}
  \begin{tabular}{lll}
    \toprule
    Dataset   & Task Type                      & Samples \\
    \midrule
    POPE \citep{liEvaluatingObjectHallucination2023}
      & Object existence probing        & 3{,}000 (3 splits $\times$ 1{,}000) \\
    MathVista \citep{luMathVistaEvaluatingMathematical2024}
      & Mathematical visual reasoning   & 1{,}000 (testmini) \\
    VQA-RAD \citep{lauDatasetClinicallyGenerated2018}
      & Medical radiology VQA           & 451 (test split) \\
    \bottomrule
  \end{tabular}
\end{table}
